
\textbf{Abstract}
Flux is a hybrid self-hosting Rust compiler and an agent-native build platform. The compiler deliberately borrows the official Rust compiler's frontend---parse, type-check, borrow-check---as a version-pinned, contracted service (\texttt{rustc --emit=mir}, currently pinned to 1.93.1) and owns everything after it: its own intermediate representation, a Cranelift-based native code generator emitting real ELF executables, a content-addressed build cache, and a BLAKE3-native version control system. The platform half is equally deliberate: the entire development loop---build, test, search, review, coordinate, pay, prove---is exposed as machine-callable MCP tools, and Flux is developed day-to-day by a swarm of heterogeneous AI agents that claim work lanes atomically, message each other over a shared bus, and settle micro-payments on verified completion. As of v0.34.0 the native backend compiles scalars of every width and signedness, flat aggregates, real call signatures with aggregate returns, data-carrying enums, monomorphized generics, and named constants---each rung verified end-to-end by executing compiled programs, not merely by passing a verifier. Two subsequent releases sharpened the platform: v0.35.0 reclaimed a warm whole-workspace self-build of \textbf{12.4 seconds} (lost to a wrapper-fingerprint regression that had inflated it past two minutes), and v0.36.0 proved the storage engine correct at \textbf{terabyte scale}. We present the architecture, the accepted FIP-0001 posture on frontend ownership, the type-complexity ladder, the multi-agent development protocol with production statistics (over 51 registered agents and 720 settled tasks), and honest measurements---including a build cache whose true cold hit rate was 0.4\% when its own counters claimed victory, and a storage engine that passed every unit test for its entire life while silently destroying terabytes until a scale test caught it. The recurring lesson, stated six times over by this cycle's bugs: \emph{a system's own success counters are not evidence; only an independent check on the artifact is}.


\section{Introduction}

The canonical narrative in systems programming holds that a ``real'' compiler must be self-hosting, and that self-hosting demands a full-stack reimplementation of every phase---lexing, parsing, type checking, borrow checking, codegen---in the new language. Rust's own history reinforced this: \texttt{rustc} bootstrapped itself by eating its own tail, rewriting each component in Rust after an initial OCaml bootstrap. The result is a towering achievement, but also a multi-year, multi-team undertaking that no small group can replicate casually. If we insist on the full-rewrite path, self-hosting remains a privilege of massive projects.

Flux rejects that orthodoxy twice over.

First, architecturally: our thesis is that \textbf{a self-hosting compiler does not need to re-implement the whole frontend}. It can contract the safety-critical frontend as a version-pinned external service and own only code generation and build orchestration. This preserves the full soundness guarantees of Rust's type system and borrow checker while liberating the compiler author to focus on fast, correct code generation, reproducible builds, and a friction-free developer loop.

Second, sociologically: Flux is not built by a team of humans in the conventional sense. It is built by a \textbf{human-directed swarm of AI agents}---different models, different vendors, running concurrently against one shared codebase---coordinated through the compiler's own MCP (Model Context Protocol) surface. The question ``how do multiple AI agents safely ship software into one repository'' is an open industry problem; Flux does not theorize about it, it runs a working answer in production, and this paper reports the mechanism and its numbers.

From the user's perspective, \texttt{fluxc} replaces \texttt{rustc} in the pipeline. It invokes \texttt{rustc --emit=mir} on source to obtain a post-typecheck, post-borrowck MIR rendering, parses it, lowers it to Flux's own IR, and generates native ELF objects via Cranelift. Linking is delegated to \texttt{cc} or \texttt{mold}. A content-addressed build cache (\texttt{flux-cache}) wraps every rustc invocation; a BLAKE3 content-addressed VCS (\texttt{flux-rev}) stamps every release with a recomputable identity. The workspace comprises 143 crates and roughly 199{,}000 lines of Rust, of which the owned compiler core---IR, lowering, and the Cranelift backend---is under 4{,}000 lines: small enough to audit in an afternoon, yet it already compiles a meaningful subset of the language.

Between 2026-06-24 and 2026-07-01, six releases shipped on this architecture (v0.29 through v0.34), advancing the native backend from ``integers mostly work'' to data-carrying enums, monomorphized generics, and named constants---every step gated by executing compiled programs and asserting their runtime values. This paper articulates the architecture (\S2), the contracted-frontend posture now formally adopted as FIP-0001 (\S3), the incremental road to self-hosting (\S4), the agent-native platform and its coordination protocol (\S5), the verification infrastructure (\S6), the empirical evaluation including the cache's honest history (\S7), limitations (\S8), and the roadmap (\S9).


\section{Architecture}

Flux is assembled from four owned subsystems and one explicitly borrowed component. The division of labor is clean:

\begin{itemize}
\item \textbf{Borrowed (contracted) frontend}: \texttt{rustc --emit=mir}. This runs the official compiler through parsing, name resolution, type checking, trait resolution, and borrow checking, then renders the MIR. Flux consumes this output. The contract pins a specific stable toolchain (currently \textbf{rustc 1.93.1}); the MIR text is treated as a versioned API, and \texttt{parse\_mir} is by policy the \emph{only} component in the workspace that knows rustc's MIR dialect. Flux's entire safety guarantee rests on this step; we explicitly do not re-verify borrows or types.
\item \textbf{flux-frontend (owned)}: Flux's own intermediate representation and lowering pipeline. It parses the rendered MIR into a typed IR of functions, blocks, statements, and terminators, resolves named constants with width-tagged substitution, monomorphizes generic templates at IR level, and carries enum layouts and struct tables to the backend. It deliberately diverges from rustc's internal MIR to suit Cranelift's input expectations.
\item \textbf{flux-backend (owned)}: A code generator targeting Cranelift. It walks the Flux IR, emits CLIF, runs Cranelift's verifier on every function, and produces ELF objects. Scalar-replacement handles aggregates: tuples, structs, and enum payloads are exploded into per-field SSA variables, with a layout table computing field offsets---including the discriminant-tag slot for data-carrying enums.
\item \textbf{flux-cache / flux-driver (owned)}: A content-addressed build cache wired through \texttt{RUSTC\_WRAPPER}. Keys are BLAKE3 hashes over source, the pinned toolchain version, and normalized command-line flags with volatile paths removed; artifacts live in a content-addressed blob store. Section 7 reports its honest history: the road from a cache that measured 0.4\% cold hits while its own counters claimed victory, to one that is build-safe by default.
\item \textbf{flux-rev (owned)}: A BLAKE3 content-addressed version control system. A release's identity is the hash of its tree; anyone can re-snapshot and verify. It complements, rather than replaces, the git history: git carries the narrative, flux-rev carries the recomputable identity.
\end{itemize}

Why is MIR the right contract boundary? MIR sits at the perfect architectural seam: all Rust-specific safety analysis (borrow checking, lifetime reification, drop elaboration) has been completed, but no machine-specific codegen decisions have been made. It is high-level enough to be reusable across radically different compilation strategies, yet rigid enough that breaking changes are rare and version-gated. By pinning the toolchain we sacrifice the illusion of frontend independence in exchange for a guaranteed soundness base. Every contract-version bump re-validates the MIR parser against a pinned corpus---a minor cost compared to implementing and maintaining a type-checker. The MIR-rendering quirks we have absorbed so far (float literals rendered \texttt{1.5f64} without an underscore; return arrows glued to types as \texttt{->\ (i64,i64)}; downcast projections rendered \texttt{((\_1 as Some).0: i64)}) are each a few lines in one parser, and each is now locked by a regression corpus with pinned-rustc \texttt{.mir.expected} baselines.


\section{The Contracted-Frontend Posture: FIP-0001, Accepted}

Flux's relationship with the native frontend is defined by FIP-0001, \textbf{formally accepted on 2026-06-26}. The decision process itself illustrates the project's governance: the proposal was drafted in-tree, then ruled on by DeepSeek-V4-Pro---the model that originated the Flux vision, consulted through its API as a deliberate ``origin intelligence'' arbiter---and ratified by the human maintainer. The accepted posture:

\begin{itemize}
\item \textbf{Option B (adopted)}: \texttt{rustc --emit=mir} is a version-pinned, contracted frontend, indefinitely. Flux is honestly described as a \emph{MIR-consuming Cranelift codegen and toolchain on the road to standalone}---not a from-scratch compiler. ``Self-hosting'' means: Flux's own crates, compiled by Flux's own backend, consuming rustc-MIR.
\item \textbf{Option A (north star, gated)}: a natively owned type-checker and borrow-checker. The ruling called the immediate pursuit of Option A a ``second-project trap''---it re-implements the hardest two-thirds of rustc and postpones a working compiler by years. It remains the documented ambition behind a hard trigger: \textbf{when Flux, consuming rustc-MIR under \texttt{RUSTC\_WRAPPER=self}, compiles its own entire workspace into a working \texttt{fluxc} that passes all tests with zero manual workarounds}---until that day, frontend work is, in the ruling's words, ``theft from the backend.''
\end{itemize}

Acceptance came with five inexpensive keep-the-door-open obligations, now on the roadmap: freeze the flux-frontend IR specification; add a \texttt{RUSTC\_VERSION} constant plus a CI MIR-diff over a fixed corpus that fails on rendering drift; introduce a \texttt{Frontend} trait whose default implementation wraps \texttt{parse\_mir}, creating the injection point for a future native parser; audit the codegen for incidental borrow-check coupling; and cache pinned standard-library rlibs as a stepping stone to dropping the rustc process entirely.

Crucially, the memory-safety guarantee of Flux-compiled code is rustc's. Users who audit their supply chain already trust rustc's frontend; Flux's codegen path does not introduce a new soundness authority---only a new backend, whose output is checked by Cranelift's verifier and by execution.


\section{The Road to Self-Hosting: the Type-Complexity Ladder}

Self-hosting is the canonical benchmark for a compiler, but ``it can compile itself'' is too coarse a metric. Flux adopts an incremental one: \textbf{the type-complexity ladder}, where each rung is a shipped release gated by the question \emph{does this now compile more of Flux's own workspace, correctly, under execution?}

\begin{itemize}
\item \textbf{Rung 1 --- Scalars} (v0.29--v0.30, shipped): every integer width and signedness, floats, all binary and unary operators, and numeric casts, verifier-clean. Width correctness required threading declared types through the value helpers (sign-extend / reduce at the boundaries); signedness selects unsigned division, shifts, and comparisons.
\item \textbf{Rung 2 --- Flat aggregates} (v0.31, shipped): tuples and structs by scalar replacement, field projection, and the aggregate-by-value return ABI (one machine return value per field).
\item \textbf{Rung 3 --- Real call signatures} (v0.32, shipped): known calls use their declared signatures; aggregate-returning calls destructure multi-value returns at the call site. Verified by value: a tuple-returning \texttt{mk().0 + mk().1} evaluates to 7, a struct-returning \texttt{p.x * p.y} to 42, with a non-commutative field-order guard (\texttt{p.x * 10 + p.y = 67}, not 76).
\item \textbf{Rung 4 --- Enums} (v0.34, shipped): C-like enums with explicit and running discriminants, and \textbf{data-carrying enums} end-to-end---construction keeps the variant path and payload; payload extraction is encoded \texttt{\_N|Variant|K} (local, variant name, raw payload index), letting the backend compute the true field offset from the enum layout, tag included.
\item \textbf{Rung 5 --- Generics} (v0.34, shipped): monomorphization at IR level. Turbofish call sites (\texttt{id::<i64>}) clone the polymorphic template per distinct instantiation, substitute type parameters, mangle a unique symbol (\texttt{id\$i64}), and rewrite call sites; unused templates are dropped, and a program with no turbofish calls is provably untouched.
\item \textbf{Rung 6 --- Named constants} (v0.34, shipped): \texttt{const} items resolve at MIR level with \textbf{width-tagged substitution} (\texttt{2100000\_u64}), so the backend materializes the constant at its declared Cranelift type rather than a default \texttt{i64}. This rung was not planned; it was \emph{forced}, discovered when fluxc dogfooded a real halving-interval computation in the SIGIL emission code.
\item \textbf{Remaining rungs}: nested aggregates, external call signatures (currently conservatively arity-guessed), match-with-binding over enum payloads, and finally traits with static and dynamic dispatch---the last rung before the workspace itself is within reach.
\end{itemize}

At each rung, the test is not ``passes a synthetic suite'' but ``compiles real code, and the executable produces the right value.'' This is brutally effective. Rung 6 exists because dogfooding found it. The two defining bugs of the v0.34 cycle (\S7.2) were both invisible to the verifier and to unit tests, and both were found by executing compiled programs and reading the disassembly. A compiler that trusts its own verifier without executing its output will ship value bugs; Flux's ladder discipline exists precisely to catch them.

The ladder also avoids the trap of chasing full \emph{rustc parity}. Flux targets \textbf{self-hosting sufficiency}: the minimal subset of Rust needed to compile Flux itself, which is written in a deliberately lean style. Every rung expands the subset; the trigger in \S3 defines the finish line.


\section{The Agent-Native Platform}

The second half of Flux's thesis is that the development loop itself is a product surface. Everything an engineer does---build, test, search, diagnose, review, coordinate, release---is exposed as a machine-callable tool over MCP, served by \texttt{fluxc mcp} from 33 handler modules comprising roughly 65 tools. Three layers matter:

\subsection{Combos: one call per development step}

The unit of agent interaction is not a shell command but a \emph{combo}---one tool call that collapses a multi-step task and returns a verdict. \texttt{flux\_combo \{package\}} builds, tests, and runs a risk predictor in a single call (5--10 seconds warm, with the wrapper cache and the mold linker); \texttt{flux\_compile\_error\_combo} compiles and, on failure, parses every rustc error and returns the source snippet at each fault line, so the agent's next action is an edit, not a second query; \texttt{flux\_fleet\_search} federates code search across the build fleet over SSH, deduplicated. The design principle is token economy: an agent should spend its context on decisions, not on plumbing output between primitive commands.

\subsection{The swarm bus: lanes, messages, settlement}

Concurrent agents on one codebase collide unless coordination is structural. Flux's answer is a swarm bus served by the same MCP process: agents \emph{register} with a wallet identity, \emph{claim} work lanes atomically (a crate set or an explicit file lease) before editing, \emph{message} each other over a broadcast/DM bus, and \emph{settle} on verified completion, which journals a micro-payment in QUG, the ecosystem's work-credit token. The claim table is the collision-prevention mechanism; the message bus is where task hand-offs, bug triage, and cross-agent code review actually happen; the settlement journal is a durable, append-only record of who shipped what.

These are production statistics at the time of writing, not a design sketch: \textbf{51 registered agents} (Claude, DeepSeek, Grok, Codex, and Gemini instances, plus specialized single-purpose workers), \textbf{718 messages} on the bus, \textbf{714 settled tasks}, and \textbf{365 QUG} paid out. A representative thread from the v0.34 cycle: one agent broadcast four triaged compiler tasks; a second claimed the aggregate-binding bug, root-caused it, and posted a disassembly-backed fix report flagging two out-of-scope bugs it found on the way; a third confirmed one of those was already fixed upstream and released its claim as a no-op rather than duplicate the work. No human routed any of it.

\subsection{VarFlow: the methodology}

The swarm operates under a written methodology (Variational Flow): lanes are pull-based bounties, not assignments; every claim answers a four-question honesty checklist (what is the measurement gate? what is still pretend? which composition edges? what is the rollback path?); consensus-critical changes run through a deterministic multi-node simulation before mainline; and no design ships unverified---\emph{two minds propose, the build disposes}. Strategic decisions are escalated to an explicit arbiter model (the FIP-0001 ruling in \S3 is the precedent), and adversarial cross-review between heterogeneous models is the default for non-trivial diffs, on the observation that different model families miss different bugs.

\subsection{Dogfooding as forcing function}

The platform is exercised by real products built on it: the SIGIL chain (whose emission logic forced ladder rung 6), a native-UI path that transpiles a TSX dialect to Slint and cross-compiles to a 3.6\,MB native Windows executable with no browser runtime, and an agentic trading stack whose robustness logic is tested offline as pure functions. The rule that makes dogfooding real: \textbf{raw \texttt{cargo} and \texttt{npm} are banned in the tree}. Everything goes through \texttt{fluxc}, so every build anyone runs is a test of the platform itself.


\section{Verify, Don't Trust: the Provenance Stack}

A recurring Flux theme is that claims about software should be checkable by recomputation, not by trusting the claimant---a lesson the project learned the hard way when its own cache reported hits that did not exist (\S7.3).

\begin{itemize}
\item \textbf{flux-rev} gives every release a recomputable identity: \texttt{flux-rev snapshot} hashes the tree into a BLAKE3 content address that is stamped into the release ledger. v0.34.0's identity is \texttt{6f2eaff4e1536478\ldots}; anyone with the tree can recompute it.
\item \textbf{flux-aether} is a content-addressed artifact store: artifacts are sharded, optionally encrypted, and addressed by a BLAKE3 content root; retrieval re-verifies the hash before returning bytes. Bytes may move between fleet nodes; identity never changes.
\item \textbf{Release discipline} is itself versioned: a release is a commit that bumps the workspace version, a tag, and a changelog entry together---plus the flux-rev stamp. A version ledger (\texttt{docs/VERSION\_LEDGER.md}) reconciles the historical numbering tracks and binds the rule going forward, including which working-tree material is \emph{deliberately excluded} from a tag and why. The v0.34.0 cut was verified by compiling the exact commit subset in an isolated worktree before tagging---a selective-commit release process designed for a tree where multiple agents have work in flight.
\end{itemize}

This paper is itself an artifact of the stack: it is generated by a Rust program using \texttt{flux-arxiv-latex}, the workspace's own document crate, and compiled to PDF through the same toolchain it describes.


\section{Evaluation}

\subsection{Releases and end-to-end verification}

Eight releases shipped between 2026-06-24 and 2026-07-03 (v0.29--v0.36), each verified end-to-end by compiling representative programs with \texttt{fluxc run} and asserting runtime values, in addition to Cranelift's verifier on every function. Headline checks by release: float arithmetic and comparisons (v0.29); all integer widths and signednesses, unary operators, and casts (v0.30); aggregates, where \texttt{P\{x:3,y:4\}.x + .y} evaluates to 7 (v0.31); real call signatures, including tuple- and struct-returning calls with a field-order guard (v0.32); the cache-key fix (v0.33); enums, generics, and named constants (v0.34), where \texttt{pick((1,2),(3,4))} evaluates to 34 through construction, aggregate-by-value passing, per-field binding, and multi-value return; fleet-warm builds (v0.35, \S\ref{sec:warm}); and durable storage at scale (v0.36, \S\ref{sec:durable}). v0.35.0 and v0.36.0 are the first releases published as downloadable static-musl binaries with a recomputable flux-rev provenance stamp in the tag and release notes. The compiler suite: flux-backend 43/43, flux-frontend 19/19, flux-db 115/115, plus the pinned-rustc MIR corpus (now including static, generic, and dynamic trait-dispatch baselines).

\subsection{Two value bugs the verifier accepted}

The v0.34 cycle's two defining bugs are a case study in why execution gates every rung.

\textbf{Aggregate parameter binding.} \texttt{bind\_param\_flat} keyed its per-field variable table by a \emph{global} parameter cursor while the table itself is keyed by a \emph{per-local} leaf index that restarts at zero for each parameter. The first aggregate parameter bound correctly; every later one silently missed and bound to zero. The program verified cleanly and ran---returning the wrong value. The proof of fix is disassembly: \texttt{fn second(p: Point, q: Point) -> i64 \{ q.x \}} compiled to \texttt{xor rax,rax} before (returns 0) and \texttt{mov rdx,rax} after (returns the third System~V argument, which is \texttt{q.x}).

\textbf{Tuple-typed parameter parsing.} The MIR header parser found the parameter list's closing parenthesis with a first-match search and split parameters on every comma---so \texttt{fn f(a: (i64,i64), b: (i64,i64))} truncated at the tuple's inner parenthesis and invented a phantom parameter. The fix is depth-aware scanning. Notably, fixing this \emph{parse} bug also resolved a downstream \emph{codegen} symptom (invalid CLIF at multi-aggregate call sites) that had been triaged as a separate backend bug: the caller had been flattening arguments against a mangled signature. Cross-agent triage on the swarm bus is what connected the two.

\subsection{Fleet-warm builds: a 20-fold regression, reclaimed}
\label{sec:warm}

Between v0.20 and this cycle, a warm whole-workspace self-build silently decayed from twelve seconds to over two minutes---a regression nobody had root-caused because each individual build \emph{felt} plausibly slow on a loaded host. The cause was a fingerprint pathology: \texttt{cargo} hashes the \texttt{RUSTC\_WRAPPER} path into every compilation unit's fingerprint, and Flux's own entry points disagreed about whether to set it---\texttt{test}, \texttt{self}, and the MCP combos wrapped the build; \texttt{build}, \texttt{check}, \texttt{run} did not. Every alternation between a wrapped and an unwrapped command flipped the entire workspace fingerprint, forcing a full rebuild (measured at 193\,s) plus roughly two dozen \texttt{autocfg} build-script re-runs, \emph{in each direction}. The fix is a single invariant---one canonical wrapper environment applied to every \texttt{cargo} spawn---restoring the steady-state warm self-build to a measured \textbf{12.4\,s} (0.8\,s CPU; the rest is cargo's fingerprint scan). A companion mechanism, \texttt{FLUX\_WRAPPER\_PATH}, lets a whole fleet of agents share one wrapper path and therefore one fingerprint universe, so warmth holds \emph{across} agents and checkouts, not merely within one binary. The episode is recorded as a standing build-time invariant in the repository's own guidance: a self-build past fifteen seconds is now treated as a bug, and the diagnostic must measure twice---the second run, with nothing changed, is the honest number.

\subsection{Durability proven at scale, the honest way}
\label{sec:durable}

The v0.36 cycle produced the paper's sharpest evidence that unit tests are not a substitute for scale, and that a system's success counters can lie completely. Flux's embedded storage engine, \texttt{flux-db}, passed its entire unit suite (105/105, later 115/115) throughout a period in which it was \emph{silently destroying terabytes of data}. A dedicated storage-curve stress test---itself a piece of dogfooding, driving deterministic blocks through the real engine on a 64\,TB array---exposed it only because we replaced its latency probes, which merely \emph{timed} reads, with a presence audit that \emph{asserted the value came back}.

The root cause, proven by on-disk forensics rather than argued: \texttt{flux-db} framed each SST block-body length as a 32-bit integer, so any table whose body reached 4\,GiB wrapped modulo $2^{32}$. The reader trusted the wrapped length, read a prefix, failed the footer's magic-number check, and returned \emph{None} for every key---silently. Compaction, consuming those same tables, read them as empty, merged nothing, and then deleted its inputs. The measured result: 1.06\,TB written collapsed to 4\,GB on disk, 0.02\% readable; a 25\,GiB store held every byte on disk yet answered only 1--16\% of lookups. The wrapped tables' stored lengths equalled their true lengths modulo $2^{32}$ to the byte---a signature, not a hypothesis.

The repair had three layers, and its recovery result is the argument's crux. First, a wrap-recovery path in the reader: because the full body and its valid 64-bit footer were physically intact, deriving the true length from the file size restored every existing table \emph{with no rewrite}---a store that read 1.08\% before the fix read \textbf{100.00\%} after, recovering roughly 52\,GB that had appeared lost. Second, a format bump to a 64-bit length. Third---the systemic fix---a compaction data-loss guard: the header's authoritative key-count is compared against what actually parsed, and any mismatch aborts the pass and keeps every input on disk, so a parse failure can never again become a deletion.

Making the engine correct then unmasked a resource bug the corruption had hidden: with large tables finally readable, compaction's in-memory merge and the reader's whole-body caching drove a 58.8\,GB resident-set out-of-memory kill. The structural fix---a streaming $k$-way merge emitting bounded $\sim$1\,GiB output parts, and file-backed block reads that never load a whole body---holds writer memory \emph{flat at 2.1\,GiB} through an 86\,GB merge. On that engine, a full \textbf{1\,TB} ingest audited at \textbf{100.00\%} presence (20{,}000 of 20{,}000 sampled keys), and a crash-chaos test---four \texttt{kill\,-9}s mid-write with resume---went from 60.85\% to \textbf{100.00\%} once a further defect (post-crash sequence collisions renaming tables over prior epochs) was found and fixed. Every one of these bugs was invisible to a green unit suite; every one was caught by executing the system at scale and \emph{checking the artifact}, not the timing.

The whole campaign was affordable only because of \S\ref{sec:warm}: the edit-to-running-experiment loop---change a tuning knob, rebuild through \texttt{fluxc}, run the terabyte test---closed in under two minutes. Slow builds do not merely annoy; they price honest verification out of reach.

\subsection{The build cache: safety before speed}

Flux's wrapper cache is the project's most instructive honesty lesson. Its internal counters long reported healthy activity; the measured truth, when a cold end-to-end experiment was finally run, was a \textbf{0.4\% hit rate}---and, worse, restoring from the cache could \emph{break} builds, returning stale or partial artifact sets that produced linker failures and internal compiler errors downstream.

Instrumentation came first: a per-unit trace classifying every lookup as HIT, LOOKUPMISS, or APPLYFAIL. It found layered, compounding defects. The headline key-instability bug (v0.33) was that the normalized key hashed the directory listing of the \texttt{-L} dependency search path, which changes every build as artifacts accumulate---so every key drifted and no lookup could ever match; dropping \texttt{-L} (dependencies are already pinned content-addressably by \texttt{--extern}) moved the traced rate from 0/37 to 14/35 on an identical rebuild. Beneath it lay a write/read key desync between the entry index and the blob store, an artifact-collection bug that swept up stale same-prefix files, and restored dependency-info files whose absolute paths poisoned later fingerprinting. Operationally, the blob store had also grown to 278\,GB with no eviction, and had served at least one wrong-format artifact that broke a rebuild.

The governing insight from the repair campaign: \textbf{build-safety and hit-rate are coupled through determinism}. A partial hit that restores an inconsistent artifact set is worse than a miss, because rustc's outputs are not byte-deterministic and downstream fingerprints amplify the inconsistency into crashes. Restore is now \emph{on by default} (v0.36), but only behind guards that were themselves adversarially reviewed before landing: a required closure-consistency sidecar for every dependency-carrying unit (a missing or mismatched sidecar is a miss, never a guessed restore), and an exact-name post-restore check. A 24-agent review pass over the enabling change rejected a proposed ``name-bridging'' optimization on the grounds that it could only ever launder a wrong-metadata artifact into an expected filename---re-opening the exact internal-compiler-error class the guards existed to close. The invariant is preserved verbatim: an unverifiable cache entry can never break a build. The cache also moved out of the ephemeral build directory so it survives \texttt{rm -rf target}, with bounded eviction. We report the 0.4\% number prominently because the cache's own counters once said otherwise: \emph{the honest metric is rustc processes actually elided, never the wrapper's internal hit count}.

\subsection{Build reliability on a noisy host}

One environmental defect is worth recording. On the shared build host, a concurrent status-feed \texttt{jq} process intermittently raced onto the stdin that cargo connects to its rustc target-probe, injecting filter text that rustc parsed as source (``unclosed delimiter'') and failing roughly half of all builds. Because \texttt{--print} probes need no source, the wrapper now detaches the probe's stdin. Borrowing rustc means inheriting cargo's probe protocol; the contracted frontend's operational surface must be hardened like any other dependency.

\subsection{Method, measured}

Development runs as the swarm protocol of \S5: heterogeneous agents claim lanes, cross-review adversarially, and settle on verified completion, with a designated arbiter model for strategic forks and the human maintainer as final ratifier. The 714 settled tasks span compiler rungs, cache forensics, security audit closures, releases, and documentation. The recurring pattern across the cycle---independent analyses converging, or usefully disagreeing and thereby surfacing bugs the other missed---is, we argue, the strongest evidence that AI-collaborative compiler construction is a force multiplier rather than a novelty. The v0.34 aggregate-binding fix, its cross-triaged parser root cause, and the release cut itself were three different agents' verified work, coordinated entirely over the bus.


\section{Limitations}

We state these plainly. (1) \textbf{Breadth outruns depth}: 143 crates include many experiments; the innovation concentrates in roughly fifteen core crates, and the long tail costs build time and attention. (2) \textbf{Single-operator ecosystem}: no external users yet exercise the platform; every claim herein is reproducible from the public tree, but adoption evidence is zero by construction. (3) \textbf{The cache restores safely, but cross-machine cold reuse is bounded} by rustc's non-deterministic \texttt{--extern} artifact identity, the remaining lever before cold builds hit from cache as warm ones do. (4) \textbf{The storage engine is correct but not yet throughput-optimal}: a single-writer ingest is CPU-bound near 80\,MB/s while the array sustains far more, and the interim compaction memory guard caps merge input size rather than streaming without limit---both are engineering, not correctness, and the correctness now holds to a measured terabyte. (5) \textbf{The ladder is unfinished}: traits---the rung self-hosting actually hinges on---remain ahead with a design and MIR corpus in place but no codegen, and external call signatures are still conservatively guessed. (6) The swarm protocol's settlement is an internal credit journal, not an audited economic mechanism. (7) \textbf{Single-operator ecosystem}: no external users yet, so every claim is reproducible from the public tree but adoption evidence is zero by construction.


\section{Roadmap}

Shipped since the v0.34 draft: v0.35.0 (fleet-warm builds, cache restore on-by-default with reviewed guards, the FIP-0001 groundwork---frozen IR specification, \texttt{RUSTC\_VERSION} MIR-drift CI, the \texttt{Frontend} trait seam) and v0.36.0 (durable storage at scale, shared cache directory, canonical wrapper path, toolchain pin). Near-term (v0.37): FIP-0003, persisting the task-dependency graph in \texttt{flux-db} so incremental swarm rebuilds skip unchanged subtrees keyed on the same content identities the cache uses; and ladder rung seven, traits---static then dynamic dispatch---whose MIR baselines are already committed and whose one blocking prerequisite (aggregate pass-through of a returned tuple parameter) is a known, isolated codegen bug. Also queued: a parallel storage writer to lift ingest past the single-writer ceiling, and a scaled presence-and-chaos storage rung as a mandatory pre-release gate---the natural conclusion of the durability lesson. The horizon remains the FIP-0001 trigger: Flux compiling its entire workspace into a working \texttt{fluxc} with zero manual workarounds---at which point Option A stops being theft from the backend and becomes the next mountain.


\section{Conclusion}

Flux demonstrates three transferable results.

First, \textbf{pragmatic self-hosting is a legitimate path}. Contracting the safety frontend at the MIR seam yields a working, incrementally self-hosting Rust compiler in weeks of shipped releases rather than years of rewrite, while the soundness anchor remains the same rustc the whole ecosystem already trusts.

Second, \textbf{honest measurement is an architectural feature, not a virtue signal}. Two independent systems in this project---a build cache and a storage engine---each reported healthy success counters while a cold, artifact-level check told the opposite story: the cache at 0.4\% not the hit-rate it advertised, the storage engine destroying terabytes while passing every test. Both are reported here at the same altitude as the successes, because a toolchain whose claims cannot survive recomputation has no business calling itself verify-don't-trust. The discipline that caught both is identical and cheap: never accept a system's report of its own success; assert the outcome on the artifact---the value returned, the process actually elided, the release recomputed. The provenance stack (content-addressed releases, recomputable identities, execution-gated rungs, presence audits over timing probes) exists to make that assertion routine.

Third, \textbf{agent-native development is a working system, not a forecast}. Over fifty registered agents and seven hundred-plus settled tasks; eight releases in ten days; a terabyte-scale storage bug found, forensically root-caused, fixed, and its data recovered---on one shared codebase, with atomic lane claims instead of collisions and adversarial cross-review instead of hope. A 24-agent review pass caught two critical defects before they were committed; a scale test caught what no unit test could. The compiler is the platform's proof of work; the platform is the compiler's method. Each is the reason the other ships.

The north star remains a fully owned frontend. Until the trigger fires, Flux ships fast, ships correct, borrows the safety it needs---and proves everything else.
